\documentclass{article}
\usepackage[left=2cm, right=2cm, top=2cm, bottom=2cm]{geometry}
\usepackage{graphicx}
\usepackage{amsmath}
\usepackage{amssymb}
\usepackage{amsthm}
\usepackage{fancyhdr}
\usepackage{verbatim}
\usepackage{listings}
\usepackage{xcolor}
\usepackage{pdfpages}
\usepackage{cancel}
\usepackage{physics}
\usepackage{esint}
\usepackage{braket}

\lstset{
    language=C++,
    basicstyle=\ttfamily\footnotesize,
    keywordstyle=\color{blue},
    commentstyle=\color{green},
    stringstyle=\color{red},
    numbers=left,
    numberstyle=\tiny\color{gray},
    frame=single,
    breaklines=true,
    captionpos=b,
}


\title{PHYS 427: Lecture \# 9}
\author{Cliff Sun}

\newtheorem{theorem}{Theorem}[section]
\newtheorem{lemma}[theorem]{Lemma}
\newtheorem{definition}[theorem]{Definition}
\newtheorem{conjecture}[theorem]{Conjecture}
\newtheorem{proposition}[theorem]{Proposition}
\newtheorem{corollary}[theorem]{Corollary}
\newtheorem{one minute paper}[theorem]{One Minute Paper}

\pagestyle{fancy}
\lhead{\textbf{\thepage}\ \ \nouppercase{\rightmark}}
\chead{PHYS 427: Lecture \# 9}
\rhead{Cliff Sun}

\begin{document}

\maketitle

\section*{Lecture Span}
\begin{enumerate}
    \item Thermodynamic potentials
\end{enumerate}

\section*{Recap}

For reversible processes, $dQ = TdS$. 

For irrversible processedm $dQ < TdS$

\section*{Thermodynamic Potentials}

Up to this point, we have covered two thermodynamic potentials:

(Fixed temperature and volume)
\begin{equation}
    F = U - TS, \quad U
\end{equation}

For some intuition, $F$ represents the maximal work you can extract from the system.   

Let $U = Q + W$, then for irrversible processes, we obtain 

\begin{equation}
    W_{on} > \Delta F
\end{equation}

Today, we will cover two more potentials. We define enthalpy:

(fixed pressure)
\begin{equation}
    H = U + pV
\end{equation}

The final energy is Gibbs Free Energy:

(temperature and pressure fixed)
\begin{equation}
    G = U + pV - TS
\end{equation}

These are the maximum amount of work you can take out of a system (reversible process). Then, 

\begin{align}
    \Delta G &= \Delta U + p \Delta V - T \Delta S \\
    &= Q + W + p \Delta V - T \Delta S \\
    &= W + p \Delta V
\end{align}

Then, $\Delta G$ would be the extra work you need to do on top of the ideal work in order to push work out of a system. 

Now, we have four quantities in the form of thermodynamic potentials:

\begin{enumerate}
    \item $U(S,V)$; $dU = TdS - pdV$ (thermodynamic identity)
    \item Helmholtz free energy: $F(T,V) = U - TS$
    $$dF = dU - TdS - S dT \iff dF = -pdV - S dT$$
    \item Enthalpy: $H(p,S) = U + pV$, then
    $$dH = dU + pdV + Vdp \iff dH = T dS + Vdp$$
    \item Gibbs Free energy: $G(p, T) = U + pV - TS$, then 
    $$dG = Vdp - S dT$$
\end{enumerate}

\section*{Maxwell's Relations}

Use mixed partials equivalences to derive the relations. Note that, 

\begin{equation}
    T = \partial_S U, \quad p = -\partial_{V} U
\end{equation}

Then, 

\begin{equation}
    \frac{\partial T}{\partial V} = -\frac{\partial p}{\partial S}
\end{equation}

Also, 

\begin{equation}
    dF = -S dT - pdV \iff -S = \partial_T F, \quad -p = \partial_V F
\end{equation}
 
Therefore, 

\begin{equation}
    \partial_{V}S  = \partial_{T}p
\end{equation}

Use $dH$ to derive 

\begin{equation}
    \partial_p T = \partial_S V
\end{equation}

Use $dG$ to derive 

\begin{equation}
    \partial_T V = -\partial_p S
\end{equation}

\subsection*{Example}

Given gas, meausre $T,V$. But we want $S$. Then, we know that the natural variables of $S$ is $U,V$.
But we want $S$ in terms of $T$ and $V$. We first note that:

\begin{equation}
    dS = \partial_T S dT + \partial_V S dV
\end{equation}

Then, introduce heat capacity. 

\begin{equation}
    C_V = \left( \frac{dQ}{dT} \right)_V = T \frac{\partial S}{\partial T}
\end{equation}

Also, 

\begin{equation}
    \frac{\partial S}{\partial V} = \frac{\partial p}{\partial T}
\end{equation}

Then, 

\begin{align}
    dS &= \partial_T S dT + \partial_V S dV \\
    &= \frac{C_V}{T}dT + \left( \partial_T p \right) dV
\end{align}

Use ideal gas law to obtain 

\begin{align}
    dS &= \frac{C_V}{T}dT + \frac{Nk_B}{V}dV\\
    S_f - S_i &= \frac{3}{2}Nk_B \ln\left( T_f/T_i \right) + Nk_B \ln\left( V_f/V_i \right)
\end{align}



\end{document}